% GENERATED FILE. Edit canonical lessons, then run npm run generate:books.
% canonical-source-hash: fnv1a64:b84090f8babe8331
% canonical-lessons: FR-C514-troupe, FR-C514-coulisses, FR-C514-decor, FR-C514-costume, FR-C514-entracte

\chapter{Au théâtre}
\label{ch:fr-au-theatre}
\hlchaptermodality{514}

\begin{chapteropening}
\textbf{By the end of this chapter:} \emph{I can talk about a theatre company and a show.}

Complete the last lesson of chapter 514: I can talk about a theatre company and a show.
\end{chapteropening}

\section[la troupe]{la troupe — (theatre) company, troupe; troop}

\label{lesson:FR-C514-troupe}



\begin{quote}
Before the new one: say the French for rehearsal, then the French for to improvise.
\end{quote}

\subsection*{You'll want to know: la troupe}
\textbf{la troupe} — \textquotedblleft{}(theatre) company, troupe; troop\textquotedblright{}.

\begin{quote}
La troupe joue la même pièce depuis deux ans.
\end{quote}

\begin{quote}
\emph{The company has been playing the same play for two years.}
\end{quote}

Also a group of soldiers or animals.

\begin{cousinweb}
\emph{la troupe} is usually traced to Frankish \emph{thorp} \textquotedblleft{}a gathering, a village\textquotedblright{} (EtymDB), a Germanic word.
\end{cousinweb}

\begin{grammarlens}[title={What you've built}]
One more word for telling, explaining and describing at B1.
\end{grammarlens}

\subsection*{Your turn}
\begin{itemize}
\raggedright
  \item \emph{Say it:} \emph{la troupe}
  \item \emph{Say it:} \emph{la troupe}, once more
  \item \emph{Say it:} a sentence of your own with \emph{la troupe}
  \item \emph{Recall:} say the French for rehearsal, then the French for to improvise, then say \emph{la troupe} again
\end{itemize}

\begin{tcolorbox}[breakable,colback=teal!4,colframe=teal!35!black,title={Before you move on}]
What does la troupe mean? (\textquotedblleft{}(theatre) company, troupe; troop\textquotedblright{}.) Say it once more.
\end{tcolorbox}
\section[les coulisses]{les coulisses — wings, backstage}

\label{lesson:FR-C514-coulisses}



\begin{quote}
Before the new one: say the French for to improvise, then the French for (theatre) company.
\end{quote}

\subsection*{You'll want to know: les coulisses}
\textbf{les coulisses} — \textquotedblleft{}wings, backstage\textquotedblright{}.

\begin{quote}
Avant le spectacle, les acteurs attendent dans les coulisses.
\end{quote}

\begin{quote}
\emph{Before the show the actors wait in the wings.}
\end{quote}

Mostly plural. \emph{Dans les coulisses} — behind the scenes.

\begin{cousinweb}
\emph{les coulisses} comes from \emph{coulisse} \textquotedblleft{}a groove that something slides in\textquotedblright{}, built on \emph{couler} \textquotedblleft{}to flow, to slide\textquotedblright{}, from Latin \emph{cōlāre} \textquotedblleft{}to strain\textquotedblright{} (Lewis \& Short): the sets slid in grooves at the side of the stage.
\end{cousinweb}

\begin{grammarlens}[title={What you've built}]
One more word for telling, explaining and describing at B1.
\end{grammarlens}

\subsection*{Your turn}
\begin{itemize}
\raggedright
  \item \emph{Say it:} \emph{les coulisses}
  \item \emph{Say it:} \emph{les coulisses}, once more
  \item \emph{Say it:} a sentence of your own with \emph{les coulisses}
  \item \emph{Recall:} say the French for to improvise, then the French for (theatre) company, then say \emph{les coulisses} again
\end{itemize}

\begin{tcolorbox}[breakable,colback=teal!4,colframe=teal!35!black,title={Before you move on}]
What does les coulisses mean? (\textquotedblleft{}Wings, backstage\textquotedblright{}.) Say it once more.
\end{tcolorbox}
\section[le décor]{le décor — set, scenery; decor}

\label{lesson:FR-C514-decor}



\begin{quote}
Before the new one: say the French for (theatre) company, then the French for wings.
\end{quote}

\subsection*{You'll want to know: le décor}
\textbf{le décor} — \textquotedblleft{}set, scenery; decor\textquotedblright{}.

\begin{quote}
Le décor représente une petite cuisine.
\end{quote}

\begin{quote}
\emph{The set shows a small kitchen.}
\end{quote}

In a theatre, often plural: \emph{les décors}.

\begin{cousinweb}
\emph{le décor} is built on \emph{décorer}, a learned borrowing from Latin \emph{decorāre} \textquotedblleft{}to adorn\textquotedblright{}, from \emph{decus} \textquotedblleft{}ornament, honour\textquotedblright{} (EtymDB; Lewis \& Short).
\end{cousinweb}

\begin{grammarlens}[title={What you've built}]
One more word for telling, explaining and describing at B1.
\end{grammarlens}

\subsection*{Your turn}
\begin{itemize}
\raggedright
  \item \emph{Say it:} \emph{le décor}
  \item \emph{Say it:} \emph{le décor}, once more
  \item \emph{Say it:} a sentence of your own with \emph{le décor}
  \item \emph{Recall:} say the French for (theatre) company, then the French for wings, then say \emph{le décor} again
\end{itemize}

\begin{tcolorbox}[breakable,colback=teal!4,colframe=teal!35!black,title={Before you move on}]
What does le décor mean? (\textquotedblleft{}Set, scenery; decor\textquotedblright{}.) Say it once more.
\end{tcolorbox}
\section[le costume]{le costume — costume; suit}

\label{lesson:FR-C514-costume}



\begin{quote}
Before the new one: say the French for wings, then the French for set.
\end{quote}

\subsection*{You'll want to know: le costume}
\textbf{le costume} — \textquotedblleft{}costume; suit\textquotedblright{}.

\begin{quote}
Les costumes de la pièce sont magnifiques.
\end{quote}

\begin{quote}
\emph{The costumes in the play are magnificent.}
\end{quote}

Also a man's suit: \emph{un costume gris}.

\begin{cousinweb}
\emph{le costume} was borrowed from Italian \emph{costume} \textquotedblleft{}custom, way of dressing\textquotedblright{} (EtymDB), which goes back to Latin \emph{cōnsuētūdō} \textquotedblleft{}habit\textquotedblright{}, from \emph{cōnsuēscere} \textquotedblleft{}to get used to\textquotedblright{} (Lewis \& Short).
\end{cousinweb}

\begin{grammarlens}[title={What you've built}]
One more word for telling, explaining and describing at B1.
\end{grammarlens}

\subsection*{Your turn}
\begin{itemize}
\raggedright
  \item \emph{Say it:} \emph{le costume}
  \item \emph{Say it:} \emph{le costume}, once more
  \item \emph{Say it:} a sentence of your own with \emph{le costume}
  \item \emph{Recall:} say the French for wings, then the French for set, then say \emph{le costume} again
\end{itemize}

\begin{tcolorbox}[breakable,colback=teal!4,colframe=teal!35!black,title={Before you move on}]
What does le costume mean? (\textquotedblleft{}Costume; suit\textquotedblright{}.) Say it once more.
\end{tcolorbox}
\section[l'entracte]{l'entracte — interval, intermission}

\label{lesson:FR-C514-entracte}



\begin{quote}
Before the new one: say the French for set, then the French for costume.
\end{quote}

\subsection*{You'll want to know: l'entracte}
\textbf{l'entracte} — \textquotedblleft{}interval, intermission\textquotedblright{}.

\begin{quote}
Pendant l'entracte, nous avons bu un café.
\end{quote}

\begin{quote}
\emph{During the interval we had a coffee.}
\end{quote}

Masculine. Literally between the acts.

\begin{cousinweb}
\emph{l'entracte} is built from \emph{entre} \textquotedblleft{}between\textquotedblright{}, from Latin \emph{inter}, and \emph{acte} \textquotedblleft{}act\textquotedblright{}, a learned borrowing from Latin \emph{āctus}, from \emph{agere} \textquotedblleft{}to do, to drive\textquotedblright{} (Lewis \& Short).
\end{cousinweb}

\begin{grammarlens}[title={What you've built}]
One more word for telling, explaining and describing at B1.
\end{grammarlens}

\subsection*{Your turn}
\begin{itemize}
\raggedright
  \item \emph{Say it:} \emph{l'entracte}
  \item \emph{Say it:} \emph{l'entracte}, once more
  \item \emph{Say it:} a sentence of your own with \emph{l'entracte}
  \item \emph{Recall:} say the French for set, then the French for costume, then say \emph{l'entracte} again
\end{itemize}

\begin{tcolorbox}[breakable,colback=teal!4,colframe=teal!35!black,title={Before you move on}]
What does l'entracte mean? (\textquotedblleft{}Interval, intermission\textquotedblright{}.) Say it once more.
\end{tcolorbox}
